\section{La convicción y la práctica del código abierto}

\subsection{Por qué insiste en el código abierto}

Chen Tianqi (陈天奇) es un firme creyente del open source y la open science. Considera que el valor del código abierto se manifiesta en varios niveles:

\begin{enumerate}
    \item \textbf{Intercambio comunitario}: compartir mutuamente, conocer a personas distintas, obtener la retroalimentación positiva de que le den «me gusta» a tu código
    \item \textbf{Una vía para que la universidad llegue a la práctica}: a la universidad misma no le resulta fácil llevar la tecnología directamente a la práctica; el código abierto es la mejor manera de compartirla con todos y construirla en conjunto
    \item \textbf{Formación de estudiantes}: en una empresa, uno normalmente empieza por el componente pequeño de un proyecto grande, pero un doctorante que hace un proyecto de código abierto tiene la oportunidad de \textbf{diseñar la arquitectura de un sistema desde cero}------ese tipo de entrenamiento es sumamente valioso
    \item \textbf{Una ventaja natural de la universidad}: neutralidad más espíritu de entrega a la comunidad
\end{enumerate}

\subsection{La clave para construir una comunidad de código abierto}

\begin{importantbox}{Los tres elementos de un proyecto de código abierto exitoso}
\begin{enumerate}
    \item \textbf{La tecnología correcta}: que todos believe in it. La tecnología misma es la base; sin una buena tecnología, la comunidad no se congrega.
    \item \textbf{La passion de mantenerlo de manera continua}: el éxito de un software no lo determina la forma que tenga su código en un momento dado, sino que necesita mantenerse constantemente. Los issues y la retroalimentación de los usuarios son «sumamente valiosos»------son el mejor insumo para la mejora iterativa.
    \item \textbf{Build trust}: si uno hace algo y de repente deja de mantenerlo, cuando lance el siguiente proyecto la gente dudará de «si en verdad hay que creer en esto». La confianza requiere acumularse a largo plazo y, una vez perdida, es muy difícil reconstruirla.
\end{enumerate}
\end{importantbox}

Chen Tianqi mismo suele «revisar con frecuencia» los issues y la retroalimentación de los usuarios. Considera que el código abierto no es free------requiere mucho tiempo para revisar la retroalimentación, corregir errores e interactuar con la comunidad, cosas que en el sistema tradicional de evaluación de la investigación no tienen retribución directa (no cuentan como publicaciones). «Así que no es que todos los grupos deban tomar este camino; casi ningún grupo lo hará, pero es una manera de hacerlo.»

El presentador preguntó de nuevo: ¿revisas uno por uno todos los issues de todos tus proyectos de código abierto? Chen Tianqi respondió que sí los revisa con frecuencia, porque «un software no puede tener su éxito determinado por la forma que tenga su código en un momento dado------necesita mantenerse constantemente».

\subsection{Resumen de la sección}

Para Chen Tianqi, el código abierto no es solo una forma de difundir tecnología, sino también un valor: que la tecnología la use más gente, que cada quien pueda crear su propio sistema de IA. Esta también fue la intención original detrás de OctoML------«para que la comunidad de código abierto pueda prosperar».
